% problem-000420 4 漂白粉有效氯分析
\section*{4. 漂白粉有效氯分析}
准确称取7.630g研细试样，⽤蒸馏⽔溶解，定容于1000mL容量瓶。

移取25.00mL该试样溶液⾄250mL锥形瓶中，加⼊过量的KI⽔溶液，以⾜量的1:1⼄酸溶液酸化，以 0.1076 mol L−1 的 $\mathrm { N a _ { 2 } S _ { 2 } O _ { 3 } }$ 标准溶液滴定⾄终点，消耗18.54 mL。

移取25.00mL试样溶液⾄250mL锥形瓶中，缓慢加⼊⾜量的 $3 \% \mathrm { H } _ { 2 } \mathrm { O } _ { 2 }$ ⽔溶液，搅拌⾄不再产⽣⽓泡。以 0.1008mol L−1 的 $\mathrm { A g N O _ { 3 } }$ 标准溶液滴定⾄终点，消耗20.36mL。

移取25.00mL试样溶液⾄100.0 mL容量瓶中，以蒸馏⽔稀释到刻度。移取25.00mL该稀释液⾄250mL锥形瓶中，以⾜量的 $3 \% \mathrm { H } _ { 2 } \mathrm { O } _ { 2 }$ ⽔溶液处理⾄不再产⽣⽓泡。于氨性缓冲液中以0.01988mol L−1的EDTA标准溶液滴定⾄终点，消耗24.43mL。

\subsection*{4-1}
计算该漂⽩粉中有效氯的百分含量（以 $\mathrm { C l } _ { 2 }$ 计）

\subsection*{4-2}
计算总氯百分含量

\subsection*{4-3}
计算总钙百分含量
